Let $G$ be a (possibly infinite) simple undirected graph.  For each $n\in\mathbb N$, define $h_G(n)$ to be the least integer $t$ such that every induced subgraph $H$ of $G$ with $|V(H)|=n$ can be made bipartite by deleting at most $t$ edges.
Equivalently,
\[
 h_G(n):=\max\{\tau(H):\;H\subseteq G,\ |V(H)|=n\},
\]
where $\tau(H)$ is the minimum number of edges that must be deleted from $H$ to make it bipartite.

Question: determine/estimate the growth of $h_G(n)$ as $n\to\infty$.  In particular, is it true that
\[
\frac{h_G(n)}{n}\to\infty\quad\text{for every graph $G$ with }\chi(G)=\aleph_1\ ?
\]
